\documentclass[a4paper,12pt]{article}

\usepackage[margin=1.8cm]{geometry}
\usepackage[T1]{fontenc}
\usepackage[utf8]{inputenc}
\usepackage[ngerman]{babel}
\usepackage{array}
\usepackage{longtable}
\usepackage[table]{xcolor}

\renewcommand{\arraystretch}{1.5}
\setlength{\tabcolsep}{6pt}

\definecolor{HeaderBlueDark}{RGB}{70,110,170}
\definecolor{RowBlueLight}{RGB}{235,243,255}
\definecolor{RowBlueDark}{RGB}{220,233,250}

\begin{document}
    
    \section{Personalpronomen}
        \rowcolors{2}{RowBlueLight}{RowBlueDark}
        \begin{longtable}{>{\bfseries}p{2.4cm}|*{9}{>{\centering\arraybackslash}p{1.2cm}}}
            \rowcolor{HeaderBlueDark}
            \color{white}\textbf{Kasus} &
            \color{white}\textbf{ich} &
            \color{white}\textbf{du} &
            \color{white}\textbf{er} &
            \color{white}\textbf{sie} &
            \color{white}\textbf{es} &
            \color{white}\textbf{wir} &
            \color{white}\textbf{ihr} &
            \color{white}\textbf{sie} &
            \color{white}\textbf{Sie} \\
            Nominativ & ich & du & er & sie & es & wir & ihr & sie & Sie \\
            Akkusativ & mich & dich & ihn & sie & es & uns & euch & sie & Sie \\
            Dativ & mir & dir & ihm & ihr & ihm & uns & euch & ihnen & Ihnen \\
            Genitiv & meiner & deiner & seiner & ihrer & seiner & unser & euer & ihrer & Ihrer \\
        \end{longtable}
    
    
    \section{Possessivpronomen / Possessivartikel}
        \rowcolors{2}{RowBlueLight}{RowBlueDark}
        \begin{longtable}{>{\bfseries}p{2.4cm}|*{9}{>{\centering\arraybackslash}p{1.2cm}}}
            \rowcolor{HeaderBlueDark}
            \color{white}\textbf{Kasus} &
            \color{white}\textbf{ich} &
            \color{white}\textbf{du} &
            \color{white}\textbf{er} &
            \color{white}\textbf{sie} &
            \color{white}\textbf{es} &
            \color{white}\textbf{wir} &
            \color{white}\textbf{ihr} &
            \color{white}\textbf{sie} &
            \color{white}\textbf{Sie} \\
            Nominativ & mein & dein & sein & ihr & sein & unser & euer & ihr & Ihr \\
            Akkusativ & meinen & deinen & seinen & ihren & sein & unseren & euren & ihren & Ihren \\
            Dativ & meinem & deinem & seinem & ihrem & seinem & unserem & eurem & ihrem & Ihrem \\
            Genitiv & meines & deines & seines & ihres & seines & unseres & eures & ihres & Ihres \\
        \end{longtable}
    
    
    \section{Reflexivpronomen}
        \rowcolors{2}{RowBlueLight}{RowBlueDark}
        \begin{longtable}{>{\bfseries}p{2.4cm}|*{9}{>{\centering\arraybackslash}p{1.2cm}}}
            \rowcolor{HeaderBlueDark}
            \color{white}\textbf{Kasus} &
            \color{white}\textbf{ich} &
            \color{white}\textbf{du} &
            \color{white}\textbf{er} &
            \color{white}\textbf{sie} &
            \color{white}\textbf{es} &
            \color{white}\textbf{wir} &
            \color{white}\textbf{ihr} &
            \color{white}\textbf{sie} &
            \color{white}\textbf{Sie} \\
            Akkusativ & mich & dich & sich & sich & sich & uns & euch & sich & sich \\
            Dativ & mir & dir & sich & sich & sich & uns & euch & sich & sich \\
        \end{longtable}
    
    
    \section{Demonstrativpronomen: der, die, das}
        \noindent
        Als Demonstrativpronomen werden diese Formen besonders betont. Sie haben dieselben Formen wie die entsprechenden Relativpronomen.
        
        \rowcolors{2}{RowBlueLight}{RowBlueDark}
        \begin{longtable}{>{\bfseries}p{2.4cm}|*{4}{>{\centering\arraybackslash}p{2.2cm}}}
            \rowcolor{HeaderBlueDark}
            \color{white}\textbf{Kasus} &
            \color{white}\textbf{Maskulin} &
            \color{white}\textbf{Feminin} &
            \color{white}\textbf{Neutrum} &
            \color{white}\textbf{Plural} \\
            Nominativ & der & die & das & die \\
            Akkusativ & den & die & das & die \\
            Dativ & dem & der & dem & denen \\
            Genitiv & dessen & deren / derer & dessen & deren / derer \\
        \end{longtable}
        
        \noindent
        \textbf{Hinweis:} \emph{deren} verweist normalerweise zurück oder steht vor einem Nomen. \emph{derer} wird besonders bei einer Vorausweisung verwendet, zum Beispiel: \emph{diejenigen, derer wir gedenken}.
    
    
    \section{Demonstrativpronomen: dieser}
        \rowcolors{2}{RowBlueLight}{RowBlueDark}
        \begin{longtable}{>{\bfseries}p{2.4cm}|*{4}{>{\centering\arraybackslash}p{2.2cm}}}
            \rowcolor{HeaderBlueDark}
            \color{white}\textbf{Kasus} &
            \color{white}\textbf{Maskulin} &
            \color{white}\textbf{Feminin} &
            \color{white}\textbf{Neutrum} &
            \color{white}\textbf{Plural} \\
            Nominativ & dieser & diese & dieses / dies & diese \\
            Akkusativ & diesen & diese & dieses / dies & diese \\
            Dativ & diesem & dieser & diesem & diesen \\
            Genitiv & dieses & dieser & dieses & dieser \\
        \end{longtable}
    
    
    \section{Demonstrativpronomen: jener}
        \rowcolors{2}{RowBlueLight}{RowBlueDark}
        \begin{longtable}{>{\bfseries}p{2.4cm}|*{4}{>{\centering\arraybackslash}p{2.2cm}}}
            \rowcolor{HeaderBlueDark}
            \color{white}\textbf{Kasus} &
            \color{white}\textbf{Maskulin} &
            \color{white}\textbf{Feminin} &
            \color{white}\textbf{Neutrum} &
            \color{white}\textbf{Plural} \\
            Nominativ & jener & jene & jenes & jene \\
            Akkusativ & jenen & jene & jenes & jene \\
            Dativ & jenem & jener & jenem & jenen \\
            Genitiv & jenes & jener & jenes & jener \\
        \end{longtable}
    
    
    \section{Demonstrativpronomen: derjenige}
        \rowcolors{2}{RowBlueLight}{RowBlueDark}
        \begin{longtable}{>{\bfseries}p{2.4cm}|*{4}{>{\centering\arraybackslash}p{2.2cm}}}
            \rowcolor{HeaderBlueDark}
            \color{white}\textbf{Kasus} &
            \color{white}\textbf{Maskulin} &
            \color{white}\textbf{Feminin} &
            \color{white}\textbf{Neutrum} &
            \color{white}\textbf{Plural} \\
            Nominativ & derjenige & diejenige & dasjenige & diejenigen \\
            Akkusativ & denjenigen & diejenige & dasjenige & diejenigen \\
            Dativ & demjenigen & derjenigen & demjenigen & denjenigen \\
            Genitiv & desjenigen & derjenigen & desjenigen & derjenigen \\
        \end{longtable}
    
    
    \section{Demonstrativpronomen: derselbe}
        \rowcolors{2}{RowBlueLight}{RowBlueDark}
        \begin{longtable}{>{\bfseries}p{2.4cm}|*{4}{>{\centering\arraybackslash}p{2.2cm}}}
            \rowcolor{HeaderBlueDark}
            \color{white}\textbf{Kasus} &
            \color{white}\textbf{Maskulin} &
            \color{white}\textbf{Feminin} &
            \color{white}\textbf{Neutrum} &
            \color{white}\textbf{Plural} \\
            Nominativ & derselbe & dieselbe & dasselbe & dieselben \\
            Akkusativ & denselben & dieselbe & dasselbe & dieselben \\
            Dativ & demselben & derselben & demselben & denselben \\
            Genitiv & desselben & derselben & desselben & derselben \\
        \end{longtable}
    
    
    \section{Demonstrativpronomen: solcher}
        \noindent
        Diese Tabelle zeigt die starke Deklination von \emph{solcher}, wenn kein Artikel davorsteht.
        
        \rowcolors{2}{RowBlueLight}{RowBlueDark}
        \begin{longtable}{>{\bfseries}p{2.4cm}|*{4}{>{\centering\arraybackslash}p{2.2cm}}}
            \rowcolor{HeaderBlueDark}
            \color{white}\textbf{Kasus} &
            \color{white}\textbf{Maskulin} &
            \color{white}\textbf{Feminin} &
            \color{white}\textbf{Neutrum} &
            \color{white}\textbf{Plural} \\
            Nominativ & solcher & solche & solches & solche \\
            Akkusativ & solchen & solche & solches & solche \\
            Dativ & solchem & solcher & solchem & solchen \\
            Genitiv & solchen & solcher & solchen & solcher \\
        \end{longtable}
    
    
    \section{Demonstrativische Verbindung: ein solcher}
        \noindent
        Nach dem unbestimmten Artikel wird \emph{solcher} wie ein Adjektiv dekliniert. Im Plural gibt es keinen unbestimmten Artikel.
        
        \rowcolors{2}{RowBlueLight}{RowBlueDark}
        \begin{longtable}{>{\bfseries}p{2.4cm}|*{4}{>{\centering\arraybackslash}p{2.2cm}}}
            \rowcolor{HeaderBlueDark}
            \color{white}\textbf{Kasus} &
            \color{white}\textbf{Maskulin} &
            \color{white}\textbf{Feminin} &
            \color{white}\textbf{Neutrum} &
            \color{white}\textbf{Plural} \\
            Nominativ & ein solcher & eine solche & ein solches & solche \\
            Akkusativ & einen solchen & eine solche & ein solches & solche \\
            Dativ & einem solchen & einer solchen & einem solchen & solchen \\
            Genitiv & eines solchen & einer solchen & eines solchen & solcher \\
        \end{longtable}
    
    
    \section{Verstärkte Demonstrativpronomen: ebendieser und ebenjener}
        \noindent
        \emph{ebendieser} wird wie \emph{dieser} dekliniert; \emph{ebenjener} wird wie \emph{jener} dekliniert.
        
        \rowcolors{2}{RowBlueLight}{RowBlueDark}
        \begin{longtable}{>{\bfseries}p{2.4cm}|*{4}{>{\centering\arraybackslash}p{2.2cm}}}
            \rowcolor{HeaderBlueDark}
            \color{white}\textbf{Kasus} &
            \color{white}\textbf{Maskulin} &
            \color{white}\textbf{Feminin} &
            \color{white}\textbf{Neutrum} &
            \color{white}\textbf{Plural} \\
            Nominativ & ebendieser / ebenjener & ebendiese / ebenjene & ebendieses / ebenjenes & ebendiese / ebenjene \\
            Akkusativ & ebendiesen / ebenjenen & ebendiese / ebenjene & ebendieses / ebenjenes & ebendiese / ebenjene \\
            Dativ & ebendiesem / ebenjenem & ebendieser / ebenjener & ebendiesem / ebenjenem & ebendiesen / ebenjenen \\
            Genitiv & ebendieses / ebenjenes & ebendieser / ebenjener & ebendieses / ebenjenes & ebendieser / ebenjener \\
        \end{longtable}
        
        \noindent
        \textbf{Zusätzlicher Hinweis:} \emph{selbst} und \emph{selber} sind unveränderlich und dienen der Hervorhebung. In der modernen Grammatik werden sie meistens nicht als eigene deklinierbare Demonstrativpronomen behandelt. Auch ein besonders betontes \emph{er}, \emph{sie} oder \emph{es} kann demonstrativ wirken, bleibt grammatisch aber ein Personalpronomen.
    
    
    \section{Relativpronomen}
        \rowcolors{2}{RowBlueLight}{RowBlueDark}
        \begin{longtable}{>{\bfseries}p{2.4cm}|*{4}{>{\centering\arraybackslash}p{2.2cm}}}
            \rowcolor{HeaderBlueDark}
            \color{white}\textbf{Kasus} &
            \color{white}\textbf{Maskulin} &
            \color{white}\textbf{Feminin} &
            \color{white}\textbf{Neutrum} &
            \color{white}\textbf{Plural} \\
            Nominativ & der & die & das & die \\
            Akkusativ & den & die & das & die \\
            Dativ & dem & der & dem & denen \\
            Genitiv & dessen & deren & dessen & deren \\
        \end{longtable}
    
    
    \section{Fragepronomen: welcher}
        \rowcolors{2}{RowBlueLight}{RowBlueDark}
        \begin{longtable}{>{\bfseries}p{2.4cm}|*{4}{>{\centering\arraybackslash}p{2.2cm}}}
            \rowcolor{HeaderBlueDark}
            \color{white}\textbf{Kasus} &
            \color{white}\textbf{Maskulin} &
            \color{white}\textbf{Feminin} &
            \color{white}\textbf{Neutrum} &
            \color{white}\textbf{Plural} \\
            Nominativ & welcher & welche & welches & welche \\
            Akkusativ & welchen & welche & welches & welche \\
            Dativ & welchem & welcher & welchem & welchen \\
            Genitiv & welches & welcher & welches & welcher \\
        \end{longtable}

\end{document}
